
\subsubsection{Overview}

The methodology has three components: (1) data collection from standardized leaderboards, (2) computation of ranking correlation statistics, and (3) bootstrap inference for uncertainty quantification.

\subsubsection{Data Collection}

\textbf{Source.} Model accuracy data are collected from Papers With Code leaderboards for ImageNet and ImageNet-V2. Papers With Code provides community-curated, standardized results using consistent evaluation protocols.

\textbf{Inclusion Criteria.} All models with reported results on both ImageNet (ILSVRC 2012 validation set) and ImageNet-V2 (MatchedFrequency variant) are included. This yielded 96 models spanning publication years 2015-2024 and architecture families including ResNet, ViT, ConvNeXt, EfficientNet, Swin, DeiT, and others.

\textbf{Rationale.} Using leaderboard data ensures comparable evaluation protocols across models. Selection bias may exist (models with poor V2 results may be underreported), but this bias likely understates ranking instability, making the finding of high stability conservative.

\subsubsection{Ranking Computation}

For each benchmark, models are ranked by top-1 accuracy in descending order (higher accuracy = better rank). Ties are handled using standard ranking conventions.

\subsubsection{Kendall-$\tau$ Correlation}

Kendall-$\tau$ (tau-b) is computed between the two ranking vectors. Kendall-$\tau$ measures the proportion of concordant vs. discordant pairs and is preferred for ordinal rankings because it is robust to ties, has a natural interpretation as proportion of correctly ordered pairs, and is less sensitive to outliers than Pearson correlation.

\textbf{Threshold.} $\tau$ < 0.90 is pre-registered as indicating significant ranking shift.

\subsubsection{Bootstrap Confidence Intervals}

95\% confidence intervals are estimated via bootstrap resampling (10,000 iterations) using the percentile method.

\subsubsection{Hypothesis Testing}

The pre-registered hypothesis is:

\begin{itemize}
\item \textbf{H1:} $\tau$ < 0.90 (rankings shift significantly)
\item \textbf{H0:} $\tau \geq$ 0.90 (rankings are preserved)
\end{itemize}

If the 95\% CI upper bound falls below 0.90, H0 is rejected in favor of H1. If $\tau \geq$ 0.90 and the CI is above 0.90, H0 is not rejected and rankings are concluded to be preserved.

